% ============================================================
% Section 3: Threat Model
% ============================================================
% Target length: ~0.5 column (300 words)
% Write: Month 3 (R2 — security person)
% This section should be written early and not change much
% ============================================================

\section{Threat Model}
\label{sec:threat_model}

\subsection{Assets}

We protect the following assets on behalf of the browser user:

\begin{itemize}
  \item \textbf{Private data} — cookies, localStorage, session
    tokens, form contents, and any data accessible via the DOM
    of pages the user has authenticated to.
  \item \textbf{System resources} — the local filesystem,
    clipboard, notification system, and download directory.
  \item \textbf{User intent} — actions performed by agents or
    extensions must reflect what the user authorised, not
    what an attacker injected.
\end{itemize}

\subsection{Adversary Model}

We consider the following adversary:

\begin{itemize}
  \item \textbf{Web content adversary.} Malicious web pages
    that embed hidden instructions (prompt injection) designed
    to manipulate an LLM agent into performing unauthorized
    actions. The adversary controls the content of visited pages
    but does not control the browser process or OS.

  \item \textbf{Malicious extension.} A browser extension
    that attempts to exceed its declared capabilities, access
    capabilities of other extensions, or exfiltrate data
    without authorization.

  \item \textbf{Compromised agent.} An LLM agent that has
    been successfully prompt-injected and issues capability
    requests beyond what the user intended. We do not assume
    the LLM itself is trustworthy.
\end{itemize}

\subsection{Trust Assumptions}

We assume the following are trusted:
the browser binary, the OS kernel, and the user's direct input
to the browser chrome (address bar, consent dialogs).
We explicitly do \emph{not} trust: web content, LLM outputs,
extension code, or the rendering engine.

\subsection{Out of Scope}

Physical attacks, OS-level compromise, and attacks on the LLM
API provider are out of scope.
Side-channel attacks on the capability broker are not addressed
in this work.

\todo{Review threat model against the four 2025 attacks to ensure
      each attack falls within scope and the adversary model
      covers the attacker's position in each case.}
